\documentclass[12pt]{article}
\usepackage[utf8]{inputenc}
\usepackage[T1]{fontenc}
\usepackage{hyperref}
\title{Woccon Language Resources for K-12 Curriculum Development}
\author{Compiled by Claude Sonnet 5.5 (claude-sonnet-5-5) for Andr\'as Kornai}
\date{}
\begin{document}
\maketitle

\section{Introduction}
Woccon (ISO 639-3 xwc; Glottocode wocc1242) is an extinct Catawban language of the Siouan family, from the Carolina coast. It is attested only in a vocabulary of about 140--150 words that John Lawson recorded and printed in 1709. No native speakers survive. Any curriculum would be a heritage and revitalization one built on this small word list.

\section{Primary Educational Resources}
No school materials were found.

\section{Linguistic and Reference Materials}
\begin{itemize}
\item Wikipedia, Woccon language: \url{https://en.wikipedia.org/wiki/Woccon_language}
\item Indian Languages Database entry, John Carter Brown Library: \url{https://www.brown.edu/Facilities/John_Carter_Brown_Library/exhibitions/ildb/details.php?id=273}
\end{itemize}
Rudes (2000), ``Resurrecting coastal Catawban'' (\emph{Southern Journal of Linguistics}), is cited in the search summary. I did not find a URL for it.

\section{Religious and Community Resources}
\begin{itemize}
\item North Carolina Department of Natural and Cultural Resources, Waccamaw Siouan Indian Tribe: \url{https://www.dncr.nc.gov/blog/2023/12/12/waccamaw-siouan-indian-tribe-d-123}
\end{itemize}
The Waccamaw are often linked to the Woccon, but scholars have not settled whether Woccon was their language. The search results also name the Cape Fear Band of the Skarure and Woccon Indians as a partner in the dictionary project below.

\section{Digital and Audio Resources}
\begin{itemize}
\item Woccon Living Dictionary (Living Tongues Institute): \url{https://livingtongues.org/woccon-living-dictionary/}
\item Woccon Living Dictionary, about page: \url{https://livingdictionaries.app/woccon/about}
\end{itemize}
The dictionary is built from Lawson's list. The next stage is to work with Siouan specialists on the correct pronunciations, so there is no reliable audio.

\section{Curriculum Development Context}
The evidence is very thin, so a curriculum would focus on heritage, history and cultural revitalization, using the living dictionary and Lawson's word list. It should be developed with descendant communities, and it should say plainly that the pronunciation is uncertain.
\end{document}